\nuevaReceta{Patatas a lo Pobre}{45 min}{receta:patatas_pobre}

    \begin{minipage}[t]{0.35\textwidth}
        \section*{Ingredientes}
        \begin{itemize}[leftmargin=*]
            \item 1 kg de patatas
            \item 1 Cebolla grande
            \item 1 Pimiento verde
            \item 1 Pimiento rojo
            \item 4 Dientes de ajo
            \item Aceite de oliva
            \item Sal
            \item Pimienta negra (opcional)
            \item 1 Chorrito de vinagre (opcional)
        \end{itemize}
    \end{minipage}
    \hfill
    \begin{minipage}[t]{0.6\textwidth}
        \section*{Preparación}
        \begin{enumerate}
            \item \textbf{Preparar las verduras:} Pelar las patatas y cortarlas en rodajas de aproximadamente medio centímetro. Cortar la cebolla en juliana y los pimientos en tiras.
            \item \textbf{Calentar el aceite:} Poner una cantidad generosa de aceite en una sartén amplia y añadir los dientes de ajo enteros y ligeramente machacados.
            \item \textbf{Añadir los ingredientes:} Incorporar las patatas, la cebolla y los pimientos. Salar y mezclar con cuidado.
            \item \textbf{Cocinar lentamente:} Tapar la sartén y cocinar a fuego medio-bajo durante 25-30 minutos, removiendo ocasionalmente sin romper demasiado las patatas.
            \item \textbf{Dorar:} Cuando las patatas estén tiernas, destapar y subir ligeramente el fuego durante unos minutos para que algunas zonas queden doradas.
            \item \textbf{Terminar:} Escurrir el exceso de aceite, ajustar de sal y, si se desea, añadir pimienta o unas gotas de vinagre antes de servir.
        \end{enumerate}
    \end{minipage}

    \vspace{1em}
    \begin{tcolorbox}[colback=orange!5, colframe=orange!50, title=Consejo, size=small]
        El secreto está en cocinar las patatas lentamente para que queden tiernas y melosas, y dorarlas solo al final. Sirven como tapa, guarnición o plato acompañado de huevos fritos.
    \end{tcolorbox}
\end{receta}
